\documentclass{handout}
\usepackage{handoutSetup}\usepackage{handoutShortcuts}   

\begin{document}\large
\handoutHeader

\begin{verbatimwrite}{\jobname.title}
Consumption and Labor Supply
\end{verbatimwrite}

\handoutNameMake

\begin{verbatimwrite}{./body.tex}
% To make handout version, set exam to false and AnswersTF to true
\provideboolean{exam}
\setboolean{exam}{true}
\setboolean{exam}{false}
\provideboolean{AnswersTF}
\setboolean{AnswersTF}{true}
%\setboolean{AnswersTF}{false}
\ifthenelse{\boolean{exam}}{\setboolean{AnswersTF}{true}}{}

Consider a consumer who has a utility function
\begin{equation}\begin{gathered}\begin{aligned}
   \uFunc(c_{t},\leisure_{t})
\end{aligned}\end{gathered}\end{equation}
where $\leisure_{t}$ is leisure (mnemonic: $z$ for laZiness) and $c_{t}$ is
consumption.  Normalize the maximum possible labor supply
to $1$; actual labor supply is $\labor_{t}$, so
that
\begin{equation}\begin{gathered}\begin{aligned}
  \labor_{t} + \leisure_{t} & =  1.
\end{aligned}\end{gathered}\end{equation}

The wage earned for working one unit of time is $\Wage_{t}$, and labor
income is the wage rate multiplied by the amount of labor supplied,
\begin{equation}\begin{gathered}\begin{aligned}
  y_{t} & =  \Wage_{t}\labor_{t}
\\  & =  (1-\leisure_{t})\Wage_{t}.
\end{aligned}\end{gathered}\end{equation}

Suppose the consumer has a fixed amount $x_{t}$ to spend in period $t$ on consumption and leisure,
\begin{equation}\begin{gathered}\begin{aligned}
  \label{eq:xdef}
  x_{t} & =  c_{t}+\leisure_{t} \Wage_{t},
\end{aligned}\end{gathered}\end{equation}
where $x_{t}$ can differ from income $y_{t}$ because this might be a 
single period in a multi-period problem.

%Thus, the consumer can `spend' income on either consumption or leisure.
\ifthenelse{\boolean{exam}}{\begin{enumerate}}{}
  \ifthenelse{\boolean{exam}}{\item Derive and explain the first order
    condition}{ The price of leisure is $\Wage_{t}$ (your income is lower
    by this amount for every extra unit of time you spend not working)
    and the price of consumption is 1, so the first order condition from the 
    optimal choice of leisure says that the ratio of the marginal
    utility of leisure to the marginal utility of consumption should
    be}{}
\begin{equation}\begin{gathered}\begin{aligned}
\Wage_{t} & =  \left(\frac{\uFunc^{\leisure}}{\uFunc^{c}}\right). \label{eq:focwucul}
\end{aligned}\end{gathered}\end{equation}

\ifthenelse{\boolean{AnswersTF}}{
\ifthenelse{\boolean{exam}}{% Beginning of exam-specific 
\begin{quote}{\it Answer:}

The }{% Beginning of handout-specific text 
To see this formally, note that the } % Beginning of common text
consumer's goal is to
\begin{equation}
\max_{\{c_{t},\leisure_{t}\}}~ \uFunc(c_{t},\leisure_{t})     .
        \label{eq:intramax}
\end{equation}
subject to a budget constraint
\begin{equation}\begin{gathered}\begin{aligned}
  \label{eq:lifebudget}
  c_{t} & =  x_{t}-\Wage_{t}\leisure_{t}
\end{aligned}\end{gathered}\end{equation}
so \eqref{eq:intramax} becomes
\begin{equation}
\max_{\{ \leisure_{t} \}} ~\uFunc(x_{t}-\Wage_{t}\leisure_{t},\leisure_{t})        
        \label{eq:intramax2}
\end{equation}
for which the FOC is
\begin{equation}\begin{gathered}\begin{aligned}
        - \uFunc^{c}\Wage_{t} + \uFunc^{\leisure} & =  0  
\\      \Wage_{t} & =  (\uFunc^{\leisure}/\uFunc^{c}).
\end{aligned}\end{gathered}\end{equation}

This is just the classic condition that says that the 
ratio of prices of two goods should equal the ratio of their
marginal utilities, which applies in any standard microeconomic
problem.  For a quantitative comparison of how this condition manifests
itself in the U.S.\ and Europe, see \cite{zwiebelLeisure}.

\ifthenelse{\boolean{exam}}{\end{quote}}{}
}{}


\ifthenelse{\boolean{exam}}{\medskip\medskip\medskip\medskip Henceforth, a}{Now, a}ssume there is an `outer' utility function ${f}(\bullet)$ which
depends on a Cobb-Douglas aggregate of consumption and leisure
\begin{equation}\begin{gathered}\begin{aligned}
  \uFunc(c_{t},\leisure_{t}) & =  {f}\left(c_{t}^{1-\leiShare}\leisure_{t}^{\leiShare}\right)
\end{aligned}\end{gathered}\end{equation}

\ifthenelse{\boolean{exam}}{
\item Show that 
}{

The inner function has the property that }$\leisure_{t} \Wage_{t} = c_{t} \eta $ for $\eta=\leiShare/(1-\leiShare)$, which implies utility can be written
\begin{equation}\begin{gathered}\begin{aligned}
  {f}((\Wage_{t}/\eta)^{-\leiShare}c_{t} ).
\end{aligned}\end{gathered}\end{equation}

\ifthenelse{\boolean{exam}}{To see this, note that the maximization problem is }{}


\ifthenelse{\boolean{AnswersTF}}{\begin{quote}

\begin{equation}\begin{gathered}\begin{aligned}
  \max_{\leisure_{t}} {f}\left((x_{t}-\leisure_{t}\Wage_{t})^{1-\leiShare}\leisure_{t}^{\leiShare}\right)
\end{aligned}\end{gathered}\end{equation}

\ifthenelse{\boolean{exam}}{{\it Answer:}}{}
FOC:
\begin{equation}\begin{gathered}\begin{aligned}
  \label{eq:23}
  (1-\leiShare) \Wage_{t}(x_{t}-\leisure_{t}\Wage_{t})^{-\leiShare}\leisure_{t}^{\leiShare}{f}^{\prime} & =  \leiShare (x_{t}-\leisure_{t}\Wage_{t})^{1-\leiShare}\leisure_{t}^{\leiShare-1}{f}^{\prime}
\\ \Wage_{t} \leisure_{t} & =  c_{t} \underbrace{\leiShare/(1-\leiShare)}_{\equiv \eta}
\end{aligned}\end{gathered}\end{equation}
so
\begin{equation}\begin{gathered}\begin{aligned}
  \label{eq:5}
  {f}(c_{t}^{1-\leiShare}\leisure_{t}^{\leiShare}) & =  {f}(c_{t}^{1-\leiShare}(\eta c_{t}/\Wage_{t})^{\leiShare})
\\ & =  {f}((\Wage_{t}/\eta)^{-\leiShare}c_{t})
\end{aligned}\end{gathered}\end{equation}
\end{quote}}





\ifthenelse{\boolean{exam}}{\item }{} 

Over long periods of time as wages have risen in the U.S., the proportion
of time spent working has not changed very much (an old stylized fact recently reconfirmed by \cite{rameyFrancisLeisure}).  Similarly, across 
countries with vastly different levels of per capita income, or across
people with vastly different levels of wages, the amount of 
variation in $\leisure_{t}$ is small compared to the size of the difference
in wages.  

\ifthenelse{\boolean{exam}}{Explain how t}{T}hese facts motivate the choice 
of utility function; \cite{kpr:prodn} show that other choices of utility
functions produce trends, but no such trends are evident in the data.  
\ifthenelse{\boolean{exam}}{Hint: T}{To see why the trends are produced, t}hink about a model in which
the lifetime lasts only a single period, with a lifetime budget constraint $\Wage_{t}  =  c_{t}+\leisure_{t}\Wage_{t}$.



\ifthenelse{\boolean{AnswersTF}}{
\ifthenelse{\boolean{exam}}{{\it Answer:}{}

\begin{quote} Using the hint along with the result, w}{W}e can solve for the level
of consumption over the lifetime as 
\begin{equation}\begin{gathered}\begin{aligned}
    \Wage_{t} & =  (1+\eta)c_{t}
\\  c_{t} & =  \Wage_{t}/(1+\eta)
\end{aligned}\end{gathered}\end{equation}
implying that leisure is
\begin{equation}\begin{gathered}\begin{aligned}
    \leisure_{t} & =  \eta c_{t}/\Wage_{t} 
\\ & =  \eta/(1+\eta)
\end{aligned}\end{gathered}\end{equation}
which is a constant (i.e.\ the amount of leisure does not trend up or down
with the level of wages).  Obviously this is what motivates the 
choice of an `inner' utility function that is Cobb-Douglas:  For such a function, people will
choose to spend constant proportions of their resources on
consumption and leisure as wages rise.
\ifthenelse{\boolean{exam}}{\end{quote}}
}{}




Now consider a two period lifetime version of the model in which
each period of life is characterized by a utility function 
of the same form and the lifetime optimization problem is
\begin{equation}\begin{gathered}\begin{aligned}
  \label{eq:26}
  \max ~~ \uFunc(c_{1},\leisure_{1}) + \DiscFac \uFunc(c_{2},\leisure_{2})
\end{aligned}\end{gathered}\end{equation}
subject to a lifetime budget constraint
  \begin{equation}\begin{gathered}\begin{aligned}
    \label{eq:DBC}
    c_{2} & =  (\Wage_{1}(1-\leisure_{1})-c_{1})\Rfree + (1-\leisure_{2}) \Wage_{2}
  \end{aligned}\end{gathered}\end{equation}



\begin{comment} % This part is the germ of a good midterm question
\item Show that if the `outer' utility function is of the CRRA form ${f}(\chi) = \chi^{1-\CRRA}/(1-\CRRA)$ then the FOC wrt $c_{1}$ implies that 


  \begin{equation}\begin{gathered}\begin{aligned}
    (\Wage_{1}/\eta)^{-\leiShare} {f}^{'}_{1} & =  \Rfree \DiscFac (\Wage_{2} \eta)^{-\leiShare} {f}^{'}_{2}
\\  \Wage_{1}^{-\leiShare} (c_{1}(\Wage_{1}/\eta)^{-\leiShare})^{-\CRRA} & =  \Rfree\DiscFac \Wage_{2}^{-\leiShare} (c_{2}(\Wage_{2}/\eta)^{-\leiShare})^{-\CRRA}
\\  \Wage_{1}^{-\leiShare(1-\CRRA)} c_{1}^{-\CRRA} & =  \Rfree\DiscFac \Wage_{2}^{-\leiShare(1-\CRRA)} c_{2}^{-\CRRA}
\\   c_{2}/c_{1} & =  (\Rfree\DiscFac)^{1/\CRRA} (\Wage_{2}/\Wage_{1})^{-\leiShare(1-\CRRA)/\CRRA} 
  \end{aligned}\end{gathered}\end{equation}
\end{comment}
  
From now on, assume that the `outer' utility function is ${f}(\chi)  = \log \chi$.
\ifthenelse{\boolean{exam}}{\item Show that t}{T}his implies that $c_{2}/c_{1} = \Rfree\DiscFac$.
\ifthenelse{\boolean{exam}}{To see why, start by noting that the budget 
constraint can be rewritten}{}

\ifthenelse{\boolean{AnswersTF}}{
\ifthenelse{\boolean{exam}}{\begin{quote} {\it Answer:}}{}


  \begin{equation}\begin{gathered}\begin{aligned}
   c_{2} & =  (\Wage_{1}-\overbrace{\leisure_{1}\Wage_{1}}^{=\eta c_{1}}-c_{1})\Rfree + \Wage_{2} - \overbrace{\Wage_{2}\leisure_{2}}^{\eta c_{2}}
\\     0  & =  (\Wage_{1}-(1+\eta) c_{1})R+\Wage_{2}-(1+\eta) c_{2}
\\ c_{2} & =  (\Rfree \Wage_{1}+\Wage_{2})/(1+\eta)- \Rfree c_{1}
  \end{aligned}\end{gathered}\end{equation}
%and notice that the $\eta$ terms are constant and can therefore be neglected
so the lifetime optimization problem becomes
  \begin{equation}\begin{gathered}\begin{aligned}
    \max_{c_{1}} \left\{\log c_{1}-\leiShare \log \Wage_{1} + \DiscFac \left(\log c_{2} - \leiShare \log  \Wage_{2} \right)\right\}
  \end{aligned}\end{gathered}\end{equation}
with FOC
\begin{equation}\begin{gathered}\begin{aligned}
  1/c_{1} & =  \Rfree \DiscFac / c_{2}
\\ c_{2}/c_{1} & =  \Rfree\DiscFac
.
\end{aligned}\end{gathered}\end{equation}

\ifthenelse{\boolean{exam}}{\end{quote}}}{}




\ifthenelse{\boolean{exam}}{ % No need to solve for c_{1} for the handout 
\item Use IBC to solve for the level of consumption $c_{1}$

\ifthenelse{\boolean{AnswersTF}}{\begin{quote}{\it Answer:}{}

  \begin{equation}\begin{gathered}\begin{aligned}
    \label{eq:16}
    PD\VFunc_{1}(c) & =  c_{1}(1+R^{-1}(\Rfree \DiscFac))
\\ & =  c_{1}(1+\DiscFac)
  \end{aligned}\end{gathered}\end{equation}

  \begin{equation}\begin{gathered}\begin{aligned}
    \label{eq:17}
    PD\VFunc_{1}(y) & =  \Wage_{1}(1-\leisure_{1}) + R^{-1}(\Wage_{2}(1-\leisure_{2}))
\\ & =  \Wage_{1}+R^{-1}\Wage_{2}-\eta (c_{1}+R^{-1} c_{2})
  \end{aligned}\end{gathered}\end{equation}

  \begin{equation}\begin{gathered}\begin{aligned}
    \label{eq:19}
    PD\VFunc_{1}(c) & =  PD\VFunc_{1}(y)
\\ c_{1}(1+\DiscFac)(1+\eta) & =  \Wage_{1}+R^{-1}\Wage_{2} \equiv h_{1}
\\ c_{1} & =  h_{1}/((1+\DiscFac)(1+\eta))
  \end{aligned}\end{gathered}\end{equation}
\end{quote}
}{}


\medskip\medskip\medskip}{}


Now we want to compare this to the two period lifetime model with no
labor supply decision.  In that model, the profile of consumption was unrelated
to the profile of labor income over lifetime.  (`Fisherian Separation' held).
\ifthenelse{\boolean{exam}}{\item Show that in this model, the profile
of labor supply satisfies
\begin{equation}\begin{gathered}\begin{aligned}
(1-\labor_{2})/(1-\labor_{1}) & =  \Rfree\DiscFac \Wage_{1}/\Wage_{2} \label{eq:LabSup}
\end{aligned}\end{gathered}\end{equation}
and explain why this makes sense in economic terms.
}{
In this model, the Fisherian Separation proposition is that the profile of $c$ is unrelated to profile of wages $\Wage$; however, the
lifetime profile of {\it leisure spending} is identical to the lifetime profile
of consumption spending,}
\ifthenelse{\boolean{AnswersTF}}{
\ifthenelse{\boolean{exam}}{\begin{quote}}{}
\begin{equation}\begin{gathered}\begin{aligned}
  \Wage_{2} \leisure_{2}/\Wage_{1}\leisure_{1} & =  \eta c_{2}/\eta c_{1} = \Rfree\DiscFac
\\ \leisure_{2}/\leisure_{1} & =  \Rfree\DiscFac \Wage_{1}/\Wage_{2}
\\ (1-\labor_{2})/(1-\labor_{1}) & =  \Rfree\DiscFac \Wage_{1}/\Wage_{2} \ifthenelse{\boolean{exam}}{}{\label{eq:LabSup}}
\end{aligned}\end{gathered}\end{equation}
so leisure moves in the opposite direction from wages, which means
labor supply $\labor = 1-z$ moves in the same direction as wages.  This
makes intuitive sense: You want to work harder when work pays better.
\ifthenelse{\boolean{exam}}{\end{quote}\medskip\medskip}{}
}{}



To make further progress, assume $\Rfree\DiscFac=1$ and define wage growth as
$\WGro=\Wage_{2}/\Wage_{1}=(1+\wGro)$.  Assume that young people tend to work about half of
their waking hours $\labor_{1}=(1/2)$ (remember vacations, weekends,
etc!).  

\ifthenelse{\boolean{exam}}{\item Show}{Note} that under these assumptions we can rewrite
\eqref{eq:LabSup} as
\ifthenelse{\boolean{exam}}{
\begin{equation}\begin{gathered}\begin{aligned}
\labor_{2}  & =  (2 \wGro + 1)/2(1+\wGro) \ifthenelse{\boolean{AnswersTF}}{\label{eq:LabSupNewer}}
\end{aligned}\end{gathered}\end{equation}
}{}
\ifthenelse{\boolean{AnswersTF}}{
\ifthenelse{\boolean{exam}}{{\it Answer:} \begin{quote}}
\begin{equation}\begin{gathered}\begin{aligned}
  (1-\labor_{2})\WGro & =  (1-\labor_{1})
\\ \wGro  & =  (1+\wGro) \labor_{2} - \labor_{1} 
\\ \labor_{2} & =  (\wGro + \labor_{1})/(1+\wGro) 
\\   & =  (2 \wGro + 1)/2(1+\wGro) \label{eq:LabSupNewer}
\end{aligned}\end{gathered}\end{equation}
\ifthenelse{\boolean{exam}}{\end{quote}}
}{}




Empirically, wages in the U.S.\ tend to grow between youth and middle
age by a factor of $\WGro \approx 2-4$ (depending on occupation and
education), so $\wGro \approx 1-3$, but labor supply 
is about the same for 55 year olds as for 25 year olds,
$\labor_{2} \approx \labor_{1}$.  


\ifthenelse{\boolean{exam}}{\item Suppose for analysis that $\wGro=2$.
  Discuss the consistency of the theory with this evidence.}{}
\ifthenelse{\boolean{AnswersTF}}{
  \ifthenelse{\boolean{exam}}{\begin{quote}{\it Answer:}
}{}
Suppose for analysis that $\wGro=2$.  Then \eqref{eq:LabSupNewer} becomes
\begin{equation}\begin{gathered}\begin{aligned}
 \labor_{2} & =  (5/6)
\end{aligned}\end{gathered}\end{equation}
so the theory says middle aged people work more than 
young people by $(2/6)/(3/6)=2/3$.  This is of course 
absurd - it implies that middle aged people would barely
have time to breathe because they were working so hard.

\ifthenelse{\boolean{exam}}{\end{quote}}
}{}

\ifthenelse{\boolean{exam}}{\medskip\medskip}{}

One objection to this analysis is that it assumed $\Rfree\DiscFac=1$, which
implies that consumption when young equals consumption when middle
aged.  In fact, on average consumption grows by about the same amount
as wages between youth and middle age.  So perhaps the right
assumption is $\Rfree\DiscFac/\WGro = 1$.  Under this assumption, we obviously
have $\labor_{2}=\labor_{1}$, matching the empirical fact.

However, there is predictably different wage growth across occupations
and education groups.  Write $\WGro_{i}=\WGro\Gamma_{i}$, where $\Gamma_{i}$ now
will differ for people in different occupations indexed by $i$, and plausible values
range from $\Gamma=0.5$ (manual laborers) to $\Gamma=1.5$ (doctors),
leaving the average value of $\Gamma$ across the two groups at $\Gamma=1$.
It is an empirical fact that the magnitude of variations in labor
supply across these groups is rather small, both in youth and in middle age.

\ifthenelse{\boolean{exam}}{\item Assuming $\Rfree\DiscFac/\WGro=1$, discuss
  whether the theory can now explain this fact.
}{Assuming $\Rfree\DiscFac/\WGro=1$,}
\ifthenelse{\boolean{AnswersTF}}{
\ifthenelse{\boolean{exam}}{

{\it Answer:}\begin{quote}}{}rewrite \eqref{eq:LabSup} for each occupation as 
\begin{equation}\begin{gathered}\begin{aligned}
  (1-\labor_{2})\Gamma_{i} & =  (1-\labor_{1})
\end{aligned}\end{gathered}\end{equation}

For $\Gamma_{i}=0.5$, if $\labor_{1} = 1/2$ we have 
\begin{equation}\begin{gathered}\begin{aligned}
  (1-\labor_{2})0.5 & =  1/2
\end{aligned}\end{gathered}\end{equation}
implying $\labor_{2} = 0$ - manual laborers would work zero hours.  However,
if $\Gamma=1.5$ so that 
\begin{equation}\begin{gathered}\begin{aligned}
   (1-\labor_{2}) (3/2) & =  1/2
\\ (1-\labor_{2}) 3 & =  1
\\ \labor_{2} & =  2/3
\end{aligned}\end{gathered}\end{equation}
so doctors would be working much harder when middle aged than when
young.  Thus, the theory says that if labor supplies are equal when
young (which is approximately true), they should differ drastically by
middle age (which is not remotely true).  That is, lifetime labor
supply does not seem to respond very much to predictable variation in
lifetime wages.  This is described in the literature as a ``small
intertemporal elasticity of labor supply.''

\ifthenelse{\boolean{exam}}{\end{quote}}
 }{}

\ifthenelse{\boolean{exam}}{\end{enumerate}}{}

\end{verbatimwrite}
\input ./body.tex

\input handoutBibMake.tex

\end{document}

